\documentclass[10pt]{article}
\usepackage[a4paper, margin=2cm]{geometry}

\usepackage[T1]{fontenc}
\usepackage{newtxtext}
\usepackage[utf8]{inputenc}
\usepackage[british]{babel}

\usepackage{named}

\usepackage{graphicx}
\usepackage{amsfonts}
\usepackage{amssymb}
\usepackage{amsmath}
\usepackage{amsthm}
\usepackage{algorithm2e}
\usepackage{url}
\usepackage{xcolor}

\newtheorem*{definition}{Definition}
\newtheorem*{theorem}{Theorem}
\newtheorem*{lemma}{Lemma}
\newtheorem*{corollary}{Corollary}

\newcommand{\ols}[1]{\mskip.5\thinmuskip\overline{\mskip-.5\thinmuskip {#1} \mskip-.5\thinmuskip}\mskip.5\thinmuskip} % overline short
\newcommand{\olsi}[1]{\,\overline{\!{#1}}} % overline short italic

\title{\vspace{-4ex}KG-ER$_0$ vs symmetry\\ - draft -}
\author{Enrico Franconi et al}
\date{\today}
\begin{document}
\maketitle

%\thispagestyle{empty}

%%%%%%%%%%%%%%%%%%%%%%%%%%%%%%%%%%%%%%%%%%%%%%%%%%%%%%%%%%%%%%%%%%%%%

\section{KG-ER$_0$ %{\small(KG-ER with simple keys, i.e., tree pattern keys of depth 1)}
}

%%%

\subsection*{Vocabulary}
$\mathcal{T}$ (Object type names $T$) \\
\phantom{}\hspace{1em}
partitioned into $\mathcal{E}$ (entity type names $E$) \\
\phantom{}\hspace{1em}
and $\mathcal{V}$ (value type names $V$)\\
$\mathcal{P}$ (relationship type names $P$, with associated arity $\alpha(P)\geq 2$)\\
$\mathcal{R}$ (role names $\rho$)

\noindent

%%%

\subsection*{Declarations}

\noindent
$\textsc{RelType}(P\ (T_1, \cdots , T_n))$ 
~~~~~{ 
\parbox[t]{\textwidth}{%
one for each $P\in\mathcal{P}$\\
$n=\alpha(P)$\\
\parbox[t]{\textwidth}{%
induces~ 
$\tau:\{(P,i)\mid P\in\mathcal P,\ 1\leq i\leq\alpha(P)\}\to\mathcal T$ ~~~with $\tau(P,i)=T_i$}}}

\noindent
$\textsc{RoleName}(P.i, \rho)$ 
~~~~~~~~~~~~~~~~~~~~~~~~~~{ 
\parbox[t]{\textwidth}{%
$1\leq i\leq \alpha(P)$\\
induces a bijection from $\mathcal{R}$ to 
$\displaystyle\{(P,i)\mid P\in\mathcal P,\ 1\leq i\leq\alpha(P)\}$
}}

%%%

\subsection*{Constraints}

$\textsc{Mand}(E, P.i)$
\hfill { $\tau(P,i)=E$}\\
%
$\textsc{Key}(P\ (P.i_1, \cdots , P.i_m))$ \hfill{ $m>0$ ~;~ \(1\leq i_j\leq\alpha(P)\) ~;~ \(i_j\neq i_k\) for $j\neq k$}\\
%
$\textsc{IsA}((E_1,\cdots,E_m)\ E)$ \hfill{ $m\geq 1$}\\
$\textsc{Disjoint}((E_1,\cdots,E_m)\ E)$ \hfill{ $m\geq 2$}\\
$\textsc{Cover}((E_1,\cdots,E_m)\ E)$ \hfill{ $m\geq 2$}

%%%

\subsection*{Semantics}

\noindent$
\begin{array}{ll}
E: & \text{a unary relation over OIDs,}\\
V: & \text{a unary relation over values,}\\
P: & \text{a relation of arity \(\alpha(P)\) over OIDs or values.}
\end{array}
$\\
OIDs and values are disjoint.
%%%
\\

%%%
\noindent
$\textsc{RelType}(P\ (T_1, \cdots , T_n))$

$P.1\subseteq T_1 ~~~\cdots~~~ P.n\subseteq T_n$

\noindent
$\textsc{Mand}(E, P.i)$

$E\subseteq P.i$

\noindent
$\textsc{Key}(P\ (P.i_1, \cdots , P.i_m))$

$P.i_1 \cdots P.i_m\rightarrow P.1\cdots P.{\alpha(P)}$

\noindent
$\textsc{IsA}((E_1,\cdots,E_m)\ E)$

$E_1\subseteq E ~~;~~\cdots ~~;~~E_m\subseteq E$

\noindent
$\textsc{Disjoint}((E_1,\cdots,E_m)\ E)$ 
 
$E_1\subseteq E ~~;~~\cdots ~~;~~E_m\subseteq E ~~;~~
E_i\cap E_j=\emptyset \text{ with } i\neq j$

\noindent
$\textsc{Cover}((E_1,\cdots,E_m)\ E)$

$E_1\subseteq E ~~;~~\cdots ~~;~~E_m\subseteq E ~~;~~
E\subseteq E_1 \cup\cdots\cup E_m$

\begin{definition}[ER-Satisfiability]
A KG-ER$_0$ schema \(\mathcal S\) is \emph{ER-satisfiable} if it is
logically satisfiable, \(\mathcal S\not\models E=\emptyset\) for each
\(E\in\mathcal E\), and \(\mathcal S\not\models P=\emptyset\) for each
\(P\in\mathcal P\).
\end{definition}

\begin{theorem}[KG-ER$_0$ cannot enforce non-vacuous symmetry]
\label{thm:kger-no-symmetry}
Let \(\mathcal S\) be an ER-satisfiable KG-ER\(_0\) schema, let
\(R\) be a relationship type of arity \(2\), and assume that
\(\mathcal S\) admits a model \(M\) containing an \(R\)-tuple
\((\alpha,\beta)\) with \(\alpha\neq\beta\); equivalently,
\(\mathcal S\not\models\forall x,y.\ R(x,y)\to x=y\).
Then \(\mathcal S\) does not entail that \(R\) is symmetric:
\[
\mathcal S\not\models \forall x,y.\ R(x, y) \rightarrow R(y, x).
\]
\end{theorem}

\begin{proof}
\phantom{}\\
\smallskip\noindent\textbf{Step 1 --- The non-diagonal hypothesis.}
Every binary relation contained in the diagonal satisfies
\(R(x,y)\to R(y,x)\) trivially (\(x=y\)).  If \(\mathcal S\) forced
\(R\) to be contained in the diagonal, then symmetry would be entailed
trivially, making the conclusion false.  The hypothesis
``\(\mathcal S\) admits a model \(M\) with \((\alpha,\beta)\in R^M\)
and \(\alpha\neq\beta\)'' rules out exactly this vacuous case.

\smallskip\noindent\textbf{Step 2 --- Chase test setup.}
Assume for contradiction that
\(\mathcal S\models\forall x,y.\ R(x,y)\to R(y,x)\).  We use the
standard fair restricted chase for tuple-generating constraints,
with quotienting for key-generated equalities.  The unary inclusion
parts of cover and disjointness constraints are chased as ordinary
unary inclusions; the covering condition
\(E\subseteq E_1\cup\cdots\cup E_n\) is handled by the usual
disjunctive branching, and disjointness constraints close exactly the
branches in which two disjoint entity types are forced to overlap.
With this convention, the standard
chase test for inclusion-style sentences says that the assumed
entailment is equivalent to the following: on every non-failing
branch of the chase of the initial instance
\[
I \;=\; \{R(a,b)\}\;\cup\;\{\text{type facts forced by the declaration of } R\},
\]
where \(a\) and \(b\) are distinct labelled nulls, the chase
quotient contains \(R(b,a)\).

We will exhibit one non-failing branch whose quotient does not
contain \(R(b,a)\), contradicting the assumed entailment.

\smallskip\noindent\textbf{Step 3 --- A non-failing branch via \(M\).}
Define a partial map \(h\) on the active domain of \(I\) by
\(h(a)=\alpha\) and \(h(b)=\beta\).  Since
\((\alpha,\beta)\in R^M\), the map \(h\) extends to a homomorphism
from \(I\) into \(M\) (the type facts of \(I\) hold in \(M\) because
they are forced by the declaration of \(R\)).  At every chase step
involving a genuine non-deterministic choice (cover constraints),
choose the alternative compatible with \(M\) via \(h\).  When a
mandatory inclusion introduces fresh existential witnesses, extend
\(h\) by mapping those witnesses to the corresponding witnesses in
\(M\), whose existence is guaranteed because \(M\) satisfies the
mandatory constraint.  These choices constrain only the disjunctive
alternatives; we continue to fire all applicable rules fairly, so
that the selected branch is a fair chase branch and its limit is a
model of \(\mathcal S\).
Along this branch:
\begin{itemize}
\item disjointness constraints are never violated (they hold in \(M\), and
\(h\) maps every chase term into the corresponding type in \(M\), so
no chase term ever ends up in two disjoint entity types);
\item cover branches are non-failing: for every branching part
\(E\subseteq E_1\cup\cdots\cup E_n\) of a cover constraint triggered by a chase term \(t\)
with \(h(t)\in E^M\), pick the branch index \(j\) for which
\(h(t)\in E_j^M\) (this exists since \(M\) satisfies the cover);
\item every key step that fires generates equalities \(u\equiv v\)
with \(h(u)=h(v)\), because if a key on a relation \(P\) fires on
two \(P\)-facts \(\phi_1,\phi_2\) agreeing on the determinant in
the chase, then \(h(\phi_1),h(\phi_2)\in P^M\) agree on the
determinant in \(M\); since \(M\) satisfies the key, they agree on
all positions, hence \(h(\phi_1[k])=h(\phi_2[k])\) for every \(k\).
Consequently, the generated equalities are consistent with \(h\); in
particular, no key firing can force \(a\equiv b\), because that
would require \(h(a)=h(b)\), i.e.\ \(\alpha=\beta\), contradicting
the non-diagonal hypothesis.  Hence the key chase does not fail.
\end{itemize}
The branch is therefore non-failing, and \(h\) extends to a
homomorphism from the entire chase into \(M\).

\smallskip\noindent\textbf{Step 4 --- The invariant to be maintained.}
Under the restricted chase fixed above, a mandatory inclusion rule fires on
an entity \(e\) only when no fact in the current instance already
witnesses \(e\) at the relevant position, and existential witnesses
are introduced as fresh, mutually distinct labelled nulls.

Identify all rules in KG-ER\(_0\) that can introduce a relationship
fact during the chase.  These are exactly the mandatory inclusions
\(E\subseteq P.i\) arising from \(\textsc{Mand}(E,P.i)\) constraints.
Background typing closures and IsA generate only unary type facts;
covers split branches; disjointness closes branches; and keys are
equality-generating.  None of these, except mandatories, introduces a
new relationship fact.

Each mandatory \(E\subseteq P.i\), when fired on an entity \(e\in E\)
in the current chase instance, creates a single \(P\)-fact
\[
P(w_1,\ldots,w_{i-1},e,w_{i+1},\ldots,w_{\alpha(P)}),
\]
where every \(w_j\) is fresh and distinct from all terms already in
the chase and from the other witnesses introduced by the same firing.
The firing condition is a property of \emph{the current instance}, not
of the rule: it checks whether some \(P\)-fact already witnesses
\(e\) at position \(i\), regardless of which rule created that fact
or which entity type triggered that earlier firing.

We will use the following joint invariant.  After each completed
chase step on the branch selected in Step~3:
\begin{enumerate}
\item[(I1)] all equivalence classes are singletons; equivalently, no
non-trivial key equality has been generated;
\item[(I2)] for every relationship type \(P\), any two distinct
\(P\)-facts disagree at every position:
if \(\phi\neq\psi\) are \(P\)-facts and \(1\leq j\leq\alpha(P)\),
then \(\phi[j]\neq\psi[j]\);
\item[(I3)] every \(R\)-fact is of exactly one of the following forms:
\begin{itemize}
\item the initial fact \(R(a,b)\);
\item \(R(e,w)\), generated by a mandatory firing
\(E\subseteq R.1\) on an entity \(e\), with \(w\) fresh at the time
of that firing;
\item \(R(w,e)\), generated by a mandatory firing
\(E\subseteq R.2\) on an entity \(e\), with \(w\) fresh at the time
of that firing;
\end{itemize}
and in particular no \(R\)-fact is literally \(R(b,a)\).
\end{enumerate}

The point of bundling these properties together is to avoid any
hidden circularity: we shall not use position uniqueness first and
prove absence of key equalities only later.  Instead, the same
induction proves position separation, absence of key equalities, and
absence of \(R(b,a)\) simultaneously.

\smallskip\noindent\textbf{Step 5 --- Base case of the invariant.}
Initially, the only relationship fact is \(R(a,b)\), together with
the type facts forced by the declaration of \(R\).  Hence no two
distinct facts of the same relationship type exist, so (I2) holds
vacuously.  No key equality has been generated, so (I1) holds.  The
only \(R\)-fact is the initial \(R(a,b)\), so (I3) holds, and
\(R(b,a)\) is absent because \(a\) and \(b\) are distinct labelled
nulls.

\smallskip\noindent\textbf{Step 6 --- Preservation by non-key steps.}
Consider a chase instance satisfying (I1)--(I3).

Unary type closure, IsA steps, and the unary inclusion parts of cover
and disjointness constraints add only unary facts.  Cover steps for
\(E\subseteq E_1\cup\cdots\cup E_n\) choose a branch, and disjointness
steps close branches; neither adds relationship facts.  Thus these
steps cannot disturb (I2), cannot create \(R(b,a)\), and cannot
generate equalities.

It remains to consider a mandatory firing.  Suppose a mandatory
\(E\subseteq P.i\) fires on an entity \(e\) and creates
\[
\phi=P(w_1,\ldots,w_{i-1},e,w_{i+1},\ldots,w_{\alpha(P)}),
\]
with fresh witnesses at all positions other than \(i\).  We compare
\(\phi\) with every previously existing \(P\)-fact \(\psi\).
\begin{itemize}
\item At the triggering position \(i\), the restricted firing
condition says that no existing \(P\)-fact already witnesses \(e\)
at position \(i\).  Hence \(\psi[i]\neq e=\phi[i]\).
\item At every other position \(j\neq i\), the value \(\phi[j]=w_j\)
is fresh, so it is distinct from every term occurring in \(\psi\).
Hence \(\psi[j]\neq\phi[j]\).
\end{itemize}
Therefore the new fact \(\phi\) disagrees with every old \(P\)-fact
at every position.  Old pairs of facts already satisfied (I2), so
(I2) is preserved.  Since a mandatory step does not generate
equalities, (I1) is preserved as well.

Now specialise this argument to \(P=R\).  The initial fact \(R(a,b)\)
already witnesses \(a\) at \(R.1\) and \(b\) at \(R.2\).  Hence no
mandatory can fire on \(a\) at \(R.1\), and no mandatory can fire on
\(b\) at \(R.2\).  A firing at \(R.1\) can only create a fact
\(R(e,w)\), with \(w\) fresh; even if \(e=b\), the second component
is fresh and therefore is not \(a\).  A firing at \(R.2\) can only
create a fact \(R(w,e)\), with \(w\) fresh; even if \(e=a\), the
first component is fresh and therefore is not \(b\).  Thus a
mandatory step preserves the three-form description in (I3) and
cannot create \(R(b,a)\).

\smallskip\noindent\textbf{Step 7 --- Preservation by key steps.}
A key step on a relationship \(P\) compares two \(P\)-facts
\(\phi\) and \(\psi\) that agree, in the current quotient, at every
position of the key determinant \(X\).  By (I1), the current quotient
is still the identity relation on terms.  Therefore agreement in the
quotient is literal equality of terms.

If \(\phi=\psi\), the key step is trivial and generates no new
equality.  If \(\phi\neq\psi\), then (I2) says that \(\phi\) and
\(\psi\) disagree at \emph{every} position.  Since every
\(\textsc{Key}\) declaration has a non-empty determinant, there is at
least one position \(j\in X\); at that position
\(\phi[j]\neq\psi[j]\).  Hence the determinant agreement required
for the key is impossible.

Consequently, no key step can fire non-trivially.  Thus key steps
preserve (I1), add no relationship facts, and therefore preserve
(I2) and (I3) as well.

By induction on the chase sequence, (I1)--(I3) hold on the entire
non-failing branch selected in Step~3.  In particular, all
equivalence classes remain singletons, and the chase quotient omits
\(R(b,a)\).

\smallskip\noindent\textbf{Step 8 --- Contradiction.}
The chase quotient along the non-failing branch from Step 3 does
not contain \(R(b,a)\), contradicting the assumed entailment that
forces \(R(b,a)\) on every non-failing branch.  Therefore
\(\mathcal S\not\models\forall x,y.\ R(x,y)\to R(y,x)\).
\end{proof}


\end{document}
